\documentclass[11pt]{article}
\usepackage[utf8]{inputenc}
\usepackage{amsmath, amssymb, amsthm}
\usepackage{booktabs}
\usepackage{geometry}
\geometry{margin=1in}
\usepackage{hyperref}

\title{ACCORD: High-Impact Findings and Strategic Upgrades for Hallucination-Resistant Systems}
\author{ACCORD Research Team}
\date{April 2026}

\begin{document}

\maketitle

\begin{abstract}
This document summarizes the key empirical and theoretical breakthroughs discovered during the refinement of the ACCORD framework. We highlight three "Clean Kill Shots" that differentiate ACCORD from existing consensus mechanisms and outline the structural upgrades implemented in the V4 submission.
\end{abstract}

\section{The "Confidence Collapse" Finding}
A critical failure of majority voting and self-consistency is the maintenance of high confidence even as agent errors correlate. In our experiments, we observed a phenomenon we term the \textbf{Confidence Collapse}:
\begin{itemize}
    \item \textbf{Finding:} When agent correlation $\kappa$ exceeds 0.8, standard voting remains $>90\%$ confident despite a $300\%$ increase in unanimous error rate (CHR).
    \item \textbf{ACCORD Upgrade:} By modeling the joint error distribution via a Gaussian Copula, ACCORD detects the latent dependency and "collapses" the posterior confidence from 0.9 to 0.4. This triggers the safety abstention logic, preventing the acceptance of correlated hallucinations.
\end{itemize}

\section{The Heterogeneity Dividend}
The "Heterogeneity Dividend" refers to the massive performance gain observed when using diverse model families (e.g., GPT-4o, Claude 3.5, Llama 3) compared to homogeneous pools (e.g., multiple instances of Llama 3).
\begin{itemize}
    \item \textbf{Homogeneous Pool:} Surety Rate 3.4\%, Error when Sure 12.9\%.
    \item \textbf{Heterogeneous Pool:} Surety Rate 67.0\%, Error when Sure 9.45\%.
    \item \textbf{Impact:} This provides the first principled basis for agent pool design, proving that structural diversity is a primary driver of consensus safety.
\end{itemize}

\section{"Patient Zero" Detection via CDG}
Traditional systems treat all hallucinations as independent events. ACCORD's Claim Dependency Graph (CDG) enables the identification of \textbf{Patient Zero}:
\begin{itemize}
    \item \textbf{Mechanism:} By performing belief propagation over logical entailment edges, ACCORD traces downstream errors to their root source.
    \item \textbf{Impact:} This allows for surgical error attribution. Instead of rejecting a full response, the system can identify the specific upstream claim that corrupted the reasoning chain, significantly improving system interpretability.
\end{itemize}

\section{V4 Structural Upgrades}
\begin{enumerate}
    \item \textbf{Conformal Aggregator:} Transitioned from empirical risk to a formal Conformal Risk Control (CRC) framework with a $\frac{n}{n+1}$ finite-sample correction.
    \item \textbf{Shrinkage-Regularized Copula:} Implemented shrinkage in $\Sigma$ estimation to ensure the correlation matrix is always positive semi-definite, even under low-sample regimes.
    \item \textbf{Loopy Belief Propagation:} Hardened the BP implementation to handle dense logical graphs with a max-iteration cap of 50.
\end{enumerate}

\end{document}




